\scanpage{94}
% archive foliation: 93
\uline{équivaut} à celle de
$\mathrm{cl}\,H^{2n}_!(X,F)-\mathrm{cl}\,H^{2n-1}_!(X,F)$
mod. les $=$ choses. On peut supposer $X$
[affine et] lisse sur $\mathbf F_p$, \inserted{(\deleted{\unsure{pure}}\inserted{\unsure{purement}} de dim $n$)} alors $H^{2n}_!(X,F)$
est dual de
\[
H^0(X,\mathcal Hom(F,\mathbf Q_\ell(n)))=H^0(X,\check F)(n),
\]
\inserted{donc} purement de poids $2n$. D'autre part, $H^{2n-1}_!(X,F)$
est dual de $H^1(X,\mathcal Hom(F,\mathbf Q_\ell(n)))$, et admet
une pièce de poids $2n-1$ \uline{et} des poids $2n-2$.
\margin{\uline{peut se raffiner}, \uline{en regardant une}
\uline{autre filtration} \uline{que par les poids},
\uline{savoir \unsure{le niveau|la dimension}}}

\deleted{Une analyse}
Si $F=\mathbf Q_\ell(0)$, \deleted{alors} donc
\deleted{le cas} si on s'occupe des fonctions $\zeta$
habituelles, la connaissance ``birationnelle
grossière'' de $\zeta$ est essentiellement équivalente à :
celle de $\pi_0(\bar X)$ comme Frob-ensemble,
et celle de l'Albanèse de $X$ sur $\mathbf F_p$\ldots
\margin{Modulo conjectures de Tate}

Une analyse plus serrée \deleted{\unsure{de ces faits}} (plus
fine que par les seuls \uline{poids}, mais en
utilisant aussi les niveaux) permet de
récupérer plus par la connaissance de
$\zeta(X\deleted{,F})$ \inserted{(disons, en prenant $F=\mathbf Q_\ell$)} modulo les
$\zeta(Y\deleted{,\unsure{F}})$ ($\dim Y<n$)
\deleted{poids $\unsure{\le\rho}$}
